% answers_short.tex —— 1986 年·全国卷·文科 · 简短答案速查
% 由 tools/make_short_answers.py 自 answers.tex 的 \gksolhead/\gkanswer 提取，
% 并与 1986/1986_数学真题_全国文_答案要点.md 逐题比对后生成，勿手改（改源 MD 后重跑）。

\begin{gknotice}
\textbf{使用说明}：本节只给最终答案便于快速核对，完整推导见解析版。答案为编写者独立推导，\textbf{非官方标准答案}。
\end{gknotice}

\gkpart{选择题}
\begin{center}\begin{tabular}{ccccccccccc}
\toprule
题号 & 1 & 2 & 3 & 4 & 5 & 6 & 7 & 8 & 9 & 10 \\
\midrule
答案 & B & C & D & A & D & B & B & A & D & C \\
\bottomrule
\end{tabular}\end{center}

\gkpart{解答题}
\begin{itemize}[leftmargin=2.6em]
  \item[\textbf{11.}] $x=\dfrac12$ 或 $x=-\dfrac32$。
  \item[\textbf{12.}] $0$。
  \item[\textbf{13.}] $3\pi$。
  \item[\textbf{14.}] $\dfrac25$。
  \item[\textbf{15.}] $-40$。
  \item[\textbf{16.}] $\dfrac{x^2}{4}-y^2=1$。
  \item[\textbf{17.}] $\text{\text{平面} }PAC\perp\text{\text{平面} }PBC$。
  \item[\textbf{18.}] $z=-\dfrac32+\dfrac{3\sqrt3}{2}i$，其辐角主值为 $\dfrac{2\pi}{3}$。
  \item[\textbf{19.}] $3y^2-2y-2x+1=0$，即 $\left(y-\dfrac13\right)^2=\dfrac23\left(x-\dfrac13\right)$，是开口向右的抛物线。
  \item[\textbf{20.}] $1680$。
  \item[\textbf{21.}] (1) $\tan3A=\dfrac ab$；(2) $(1+2\cos2A)^2=a^2+b^2$
  \item[\textbf{22.}] (1) $a_n=\dfrac32-\dfrac{1}{2\cdot3^n}$；(2) $\displaystyle\lim_{n\to\infty}a_n=\dfrac32$
\end{itemize}
